\section{Ces\`aro fixed points for quantum channels}

\begin{definition}[Ces\`aro mean]\label{def:cesaro_mean}
    \lean{cesaroMean}
    \leanok
    The \emph{Ces\`aro mean} of a linear map $E$ at order $N$,
    applied to a starting matrix $\rho_0$, is
    \begin{align}
        \sigma_N
         & =\frac{1}{N}\sum_{n=0}^{N-1}E^n(\rho_0).
        \label{eq:channel_cesaro_mean}
    \end{align}
\end{definition}

\begin{theorem}[Ces\`aro telescope]\label{thm:cesaro_telescope}
    \lean{cesaroMean_telescope}
    \leanok
    \uses{def:cesaro_mean}
    For $N\geq1$, the Ces\`aro mean satisfies
    \begin{align}
        E(\sigma_N)-\sigma_N
         & =\frac{1}{N}(E^N(\rho_0)-\rho_0).
        \label{eq:channel_cesaro_telescope}
    \end{align}
\end{theorem}

\begin{proof}\leanok\uses{def:cesaro_mean}
    Applying $E$ to (\ref{eq:channel_cesaro_mean}) gives
    $E(\sigma_N)=N^{-1}\sum_{n=1}^{N}E^n(\rho_0)$.
    Subtracting $\sigma_N$ cancels all but the first and last terms, yielding
    (\ref{eq:channel_cesaro_telescope}).
\end{proof}


\begin{lemma}[Dirichlet's simultaneous approximation]
    \label{lem:dirichlet_simultaneous}
    \lean{Dirichlet.exists_int_near_mul_simultaneous}\leanok
    \uses{}
    Let $x_1,\dots,x_m$ be $m$ real numbers and $q>1$ any integer.
    Then there exist integers $n,p_1,\dots,p_m$ such that
    $1\le n\le q^m$ and $|x_k n-p_k|\le 1/q$ for all $k$.
    This is \cite[Lemma~6.1]{Wolf2012Quantum}.
\end{lemma}

\begin{proof}\leanok
    Pigeonhole principle on the integer lattice of side $q^m$: the
    $q^m+1$ points $(\{nx_1\},\dots,\{nx_m\})$ for $n=0,\dots,q^m$
    lie in $q^m$ cells; two share a cell, and their difference yields $n$.
\end{proof}


\begin{corollary}[Recurrent powers of finitely many phases]
    \label{cor:dirichlet_phase_recurrence}
    \lean{Dirichlet.exists_ge_pow_sub_one_norm_le,
        Dirichlet.exists_strictMono_pow_tendsto_one}
    \leanok
    \uses{lem:dirichlet_simultaneous}
    For finitely many complex numbers $\theta_j$ satisfying $|\theta_j|=1$,
    there is a strictly increasing sequence of positive integers $n_i$ such
    that $\theta_j^{n_i}\longrightarrow1$ simultaneously for every~$j$.
\end{corollary}

\begin{proof}\leanok
    \uses{lem:dirichlet_simultaneous}
    Write each phase as $\theta_j=e^{2\pi i x_j}$. Apply simultaneous
    approximation with successively larger denominators and choose the
    approximating exponents recursively beyond the preceding one.
\end{proof}

\begin{definition}[Peripheral spectral splitting]
    \label{def:peripheral_spectral_splitting}
    \lean{Module.End.peripheralSubspace,
        Module.End.nonPeripheralSubspace}
    \leanok
    Let $T$ be an endomorphism of a finite-dimensional complex vector space.
    Its peripheral subspace is the sum of the maximal generalized eigenspaces
    for eigenvalues $\mu$ with $|\mu|=1$; its non-peripheral subspace is the
    corresponding sum for the remaining eigenvalues. This is the splitting
    induced by the full spectral decomposition of
    \cite[Equation~(6.5)]{Wolf2012Quantum}, with the peripheral projection
    itself defined in \cite[Equation~(6.12)]{Wolf2012Quantum}.
\end{definition}

\begin{theorem}[Complementarity of the peripheral splitting]
    \label{thm:peripheral_spectral_complementarity}
    \lean{Module.End.isCompl_peripheralSubspace_nonPeripheralSubspace}
    \leanok
    \uses{def:peripheral_spectral_splitting}
    The peripheral and non-peripheral spectral subspaces are complementary.
\end{theorem}

\begin{lemma}[Powers vanish on the non-peripheral subspace]
    \label{lem:non_peripheral_pow_vanishing}
    \lean{Module.End.tendsto_pow_apply_zero_of_mem_nonPeripheralSubspace}
    \leanok
    \uses{def:peripheral_spectral_splitting}
    Let $T$ be an endomorphism of a finite-dimensional complex vector space
    whose every eigenvalue has modulus at most one. Then $T^{n}(Y)\to0$ for
    every $Y$ in the non-peripheral subspace:
    \begin{align}
        T^{n}(Y)&\longrightarrow0.
        \label{eq:non_peripheral_pow_vanishing}
    \end{align}
\end{lemma}

\begin{proof}\leanok
    Every non-peripheral eigenvalue $\mu$ has modulus strictly below one.
    On the corresponding maximal generalized eigenspace write
    $T=\mu\Id+N$ with $N^{\ell}=0$; the binomial expansion gives
    \begin{align}
        (\mu\Id+N)^{n}
        &=\sum_{m=0}^{\ell-1}\binom{n}{m}\mu^{n-m}N^{m}
        \longrightarrow0,
        \label{eq:non_peripheral_jordan_decay}
    \end{align}
    since each summand carries the factor $\mu^{n}\to0$ against a polynomial
    in $n$. The non-peripheral subspace is the sum of these eigenspaces, and
    convergence to zero is preserved by finite sums.
\end{proof}

\begin{lemma}[Powers on an eigenspace]
    \label{lem:pow-on-eigenspace}
    \lean{Module.End.pow_apply_of_mem_eigenspace}
    \leanok
    Let $T$ be an endomorphism of a vector space over a field. If
    $T(x)=\mu x$, then
    \begin{align}
        T^n(x)
         & =\mu^n x
        \notag
    \end{align}
    for every $n\geq0$.
\end{lemma}

\begin{proof}\leanok
    Induct on $n$. The induction step follows from linearity:
    $T(\mu^n x)=\mu^nT(x)=\mu^{n+1}x$.
\end{proof}

\begin{lemma}[Recurrence on a finite sum of eigenspaces]
    \label{lem:finite-eigenspace-recurrence}
    \lean{Module.End.tendsto_pow_apply_self_of_mem_iSup_eigenspace}
    \leanok
    Let $T$ be an endomorphism of a complex normed space,
    let $S\subseteq\C$ be finite, and let $(n_k)_{k\geq0}$ be a sequence of
    nonnegative integers such that $\mu^{n_k}\to1$ for every $\mu\in S$.
    If $x$ belongs to the sum of the $\mu$-eigenspaces for $\mu\in S$, then
    $T^{n_k}(x)\to x$.
\end{lemma}

\begin{proof}\leanok
    \uses{lem:pow-on-eigenspace}
    Write $x=\sum_{\mu\in S}x_\mu$ with $T(x_\mu)=\mu x_\mu$. By
    Lemma~\ref{lem:pow-on-eigenspace},
    $T^{n_k}(x_\mu)=\mu^{n_k}x_\mu\to x_\mu$. The conclusion follows by
    summing over the finite set $S$.
\end{proof}

\begin{definition}[Peripheral spectral projections]
    \label{def:peripheral_spectral_projections}
    \lean{Module.End.peripheralProjection,
        Module.End.peripheralWeightedProjection}
    \leanok
    \uses{def:peripheral_spectral_splitting,
        thm:peripheral_spectral_complementarity}
    The map $T_\varphi$ is the projection onto the peripheral subspace along
    the non-peripheral subspace. The phase-weighted peripheral map is
    $T_\varphi'=T\circ T_\varphi$.
\end{definition}

\begin{theorem}[Properties of the peripheral spectral projections]
    \label{thm:peripheral_spectral_projection_properties}
    \lean{Module.End.range_peripheralProjection,
        Module.End.isIdempotentElem_peripheralProjection,
        Module.End.peripheralProjection_comp,
        Module.End.peripheralWeightedProjection_eq_comp,
        Module.End.peripheralWeightedProjection_eq_peripheralProjection_comp,
        Module.End.peripheralProjection_apply_of_mem_eigenspace,
        Module.End.peripheralWeightedProjection_apply_of_mem_eigenspace}
    \leanok
    \uses{def:peripheral_spectral_projections}
    The range of $T_\varphi$ is the peripheral subspace, $T_\varphi$ is
    idempotent, and
    \begin{align}
        T_\varphi' & =T\circ T_\varphi=T_\varphi\circ T.
    \end{align}
    On a peripheral $\mu$-eigenvector, $T_\varphi$ acts as the identity and
    $T_\varphi'$ acts as multiplication by~$\mu$. These are
    \cite[Equations~(6.12) and~(6.13)]{Wolf2012Quantum}.
\end{theorem}

\begin{theorem}[Recurrent subsequence and peripheral projection]
    \label{thm:peripheral_projection_recurrent_subsequence}
    \lean{IsPositiveMap.tendsto_pow_apply_peripheralProjection,
        IsPositiveMap.exists_strictMono_tendsto_pow_peripheralProjection,
        IsPositiveMap.tendsto_endEquiv_pow_peripheralProjection,
        IsPositiveMap.exists_strictMono_tendsto_pow_peripheralProjection_clm}
    \leanok
    \uses{cor:dirichlet_phase_recurrence,
        def:peripheral_spectral_projections,
        thm:peripheral_jordan_trivial,
        thm:positive_trace_preserving_mean_ergodic_projection}
    Let $T:\MN{D}\to\MN{D}$ be positive and trace-preserving, with $D>0$.
    There is a strictly increasing sequence of positive integers $n_i$ such
    that $T^{n_i}\to T_\varphi$ pointwise and in operator norm. This is
    \cite[Proposition~6.3(i)]{Wolf2012Quantum}.
\end{theorem}

\begin{proof}\leanok
    \uses{cor:dirichlet_phase_recurrence,
        def:peripheral_spectral_projections,
        lem:non_peripheral_pow_vanishing,
        thm:peripheral_jordan_trivial,
        thm:positive_trace_preserving_mean_ergodic_projection}
    Bounded forward orbits imply that every eigenvalue has modulus at most
    one. On a non-peripheral generalized eigenspace, write the corresponding
    Jordan block as $\mu\Id+N$, where $N^\ell=0$. Then
    \begin{align}
        (\mu\Id+N)^n
        =\sum_{m=0}^{\ell-1}\binom{n}{m}\mu^{n-m}N^m,
    \end{align}
    and every summand tends to zero because $|\mu|<1$.
    Theorem~\ref{thm:peripheral_jordan_trivial} removes nilpotent parts on the
    peripheral subspace, while the simultaneous recurrence corollary chooses a
    common subsequence on which all peripheral
    phases tend to one. The complementary splitting gives pointwise
    convergence, which is equivalent to operator-norm convergence in finite
    dimension.
\end{proof}

\begin{corollary}[Channel structure of the peripheral projections]
    \label{cor:peripheral_projection_channel}
    \lean{IsPositiveMap.peripheralProjection_isPositiveMap,
        IsPositiveMap.peripheralProjection_isTracePreservingMap,
        IsPositiveMap.peripheralWeightedProjection_isPositiveMap,
        IsPositiveMap.peripheralWeightedProjection_isTracePreservingMap,
        IsPositiveMap.peripheralProjection_isCPMap,
        IsChannel.peripheralProjection,
        IsChannel.peripheralWeightedProjection}
    \leanok
    \uses{thm:peripheral_projection_recurrent_subsequence,
        def:peripheral_spectral_projections}
    The peripheral projection of a positive trace-preserving map is positive
    and trace-preserving, and then the weighted projection $T_\varphi'$ is
    positive and trace-preserving as well. If the original map is completely
    positive, then $T_\varphi$ and $T_\varphi'$ are quantum channels. This
    is the preservation assertion in the opening clause of
    \cite[Proposition~6.3]{Wolf2012Quantum} (item~(ii) itself states only the
    composition identity).
\end{corollary}

\begin{proof}\leanok
    \uses{thm:peripheral_projection_recurrent_subsequence}
    Every power of $T$ preserves positivity and trace, and every power of a
    completely positive map is completely positive. These properties are
    closed under the operator-norm limit. Composition with $T$ gives the
    assertion for $T_\varphi'$.
\end{proof}

\begin{lemma}[Pointwise limits of endomorphisms in finite dimension]
    \label{lem:tendsto_clm_of_tendsto_apply}
    \lean{ContinuousLinearMap.tendsto_of_tendsto_apply_of_finiteDimensional}
    \leanok
    Let $E$ be a finite-dimensional complex normed space and let $S_N$ and $P$
    be continuous linear endomorphisms of $E$. If $S_N(x)\to P(x)$ for every
    $x\in E$, then $S_N\to P$ in the operator norm.
\end{lemma}

\begin{proof}\leanok
    Fix a basis $e_1,\dots,e_m$ of $E$. Evaluation on the basis,
    \begin{align}
        \Phi(S)=(S(e_1),\dots,S(e_m)),
    \end{align}
    is a linear bijection from the endomorphisms of $E$ onto $E^m$, since an
    endomorphism is determined by, and may be prescribed arbitrarily on, a
    basis. In finite dimension every linear map is continuous, so both $\Phi$
    and $\Phi^{-1}$ are continuous. The hypothesis says $\Phi(S_N)\to\Phi(P)$
    coordinatewise, hence $S_N=\Phi^{-1}(\Phi(S_N))\to\Phi^{-1}(\Phi(P))=P$.
\end{proof}

\begin{lemma}[Closedness of the completely positive maps]
    \label{lem:cp_maps_closed}
    \lean{isClosed_setOf_isCPMap, IsCPMap.of_tendsto_toCLM}
    \leanok
    \uses{def:cp_map, thm:cp_iff_choi_posSemidef}
    Let $D>0$. The completely positive maps form a closed subset of the
    endomorphisms of $\MN{D}$ in the operator norm. In particular, if every
    $S_N:\MN{D}\to\MN{D}$ is completely positive and $S_N\to P$ in the
    operator norm, then $P$ is completely positive.
\end{lemma}

\begin{proof}\leanok
    \uses{def:cp_map, def:choi_matrix, thm:cp_iff_choi_posSemidef}
    By Theorem~\ref{thm:cp_iff_choi_posSemidef} a map is completely positive
    precisely when its Choi matrix is positive semidefinite. The Choi matrix
    depends linearly, hence continuously, on the map, and the positive
    semidefinite matrices are closed. The set of completely positive maps is
    therefore the preimage of a closed set under a continuous map.
\end{proof}

\begin{corollary}[Channel structure of the mean-ergodic projection]
    \label{cor:mean_ergodic_projection_channel}
    \lean{IsCPMap.meanErgodicProjection_isCPMap,
        IsChannel.meanErgodicProjection}
    \leanok
    \uses{def:cp_map,
        def:has_bounded_orbits,
        def:mean_ergodic_projection,
        thm:mean_ergodic_bounded_orbits,
        thm:positive_trace_preserving_mean_ergodic_projection}
    Let $T:\MN{D}\to\MN{D}$ be completely positive with bounded orbits. Then
    the mean-ergodic projection $T_\infty$ of
    \cite[Equation~(6.14)]{Wolf2012Quantum} is completely positive. If $T$ is
    completely positive and trace-preserving, the bounded orbits are automatic
    and $T_\infty$ is completely positive and trace-preserving, hence a
    quantum channel. Together with
    Corollary~\ref{cor:peripheral_projection_channel} this completes the
    preservation assertion of \cite[Proposition~6.3]{Wolf2012Quantum} for all
    three maps $T_\infty$, $T_\varphi$, and $T_\varphi'$.
\end{corollary}

\begin{proof}\leanok
    \uses{def:mean_ergodic_projection,
        thm:mean_ergodic_bounded_orbits,
        thm:positive_trace_preserving_mean_ergodic_projection,
        thm:trace_preserving_mean_ergodic_projection,
        lem:tendsto_clm_of_tendsto_apply,
        lem:cp_maps_closed}
    For $D=0$ the algebra is trivial and every linear map is completely
    positive, so assume $D>0$. Write
    \begin{align}
        S_N=\frac{1}{N}\sum_{n=0}^{N-1}T^n
    \end{align}
    for the $N$-th Ces\`aro average. Powers of a completely positive map are
    completely positive, and so are finite sums and nonnegative multiples of
    completely positive maps, so every $S_N$ is completely positive. For every
    $X\in\MN{D}$, (\ref{eq:channel_mean_ergodic_limit}) gives
    \begin{align}
        S_N(X)\longrightarrow T_\infty(X),
    \end{align}
    and Lemma~\ref{lem:tendsto_clm_of_tendsto_apply} upgrades this to
    \begin{align}
        S_N\longrightarrow T_\infty
    \end{align}
    in the operator norm. Lemma~\ref{lem:cp_maps_closed} then makes $T_\infty$
    completely positive. Trace preservation of $T_\infty$ is
    Theorem~\ref{thm:trace_preserving_mean_ergodic_projection}, and the
    bounded orbits it requires come from
    Theorem~\ref{thm:positive_trace_preserving_mean_ergodic_projection}.
\end{proof}

\begin{theorem}[Peripheral absorption of the mean-ergodic projection]
    \label{thm:peripheral_mean_ergodic_absorption}
    \lean{IsPositiveMap.peripheralProjection_comp_meanErgodicProjection}
    \leanok
    \uses{thm:peripheral_projection_recurrent_subsequence,
        def:peripheral_spectral_projections,
        def:mean_ergodic_projection,
        thm:mean_ergodic_bounded_orbits,
        thm:positive_trace_preserving_mean_ergodic_projection}
    If $T_\infty$ denotes the mean-ergodic projection of a positive
    trace-preserving map, then
    \begin{align}
        T_\varphi T_\infty=T_\infty.
    \end{align}
    This is the absorption identity adjacent to
    \cite[Equation~(6.14)]{Wolf2012Quantum}.
\end{theorem}

\begin{proof}\leanok
    \uses{thm:peripheral_projection_recurrent_subsequence,
        def:mean_ergodic_projection,
        thm:mean_ergodic_bounded_orbits,
        thm:positive_trace_preserving_mean_ergodic_projection}
    The range of $T_\infty$ consists of fixed points of~$T$, hence every power
    of $T$ fixes this range. Passing to the recurrent-subsequence limit shows
    that $T_\varphi$ also fixes it pointwise.
\end{proof}

\begin{theorem}[Asymptotic image]
    \label{thm:peripheral_asymptotic_image}
    \lean{Module.End.map_eigenspace_of_ne_zero,
        IsPositiveMap.fixedPointsSubmodule_peripheralProjection,
        IsPositiveMap.range_peripheralProjection_eq_iSup_eigenspace,
        IsPositiveMap.span_stationaryDensity_peripheralProjection_eq_peripheralSubspace,
        IsPositiveMap.exists_posSemidef_span_eq_iSup_eigenspace,
        IsPositiveMap.map_peripheralSubspace,
        IsPositiveMap.asymptotic_image}
    \leanok
    \uses{def:peripheral_spectral_projections}
    Let $T:\MN{D}\to\MN{D}$ be positive and trace-preserving, with $D>0$, and
    let
    \begin{align}
        \mathcal X_T
        =\spn\{X\in\MN{D}\mid
        \exists\varphi\in\R:\ T(X)=e^{i\varphi}X\}
    \end{align}
    be the span of its peripheral eigenvectors. Then
    \begin{enumerate}
        \item $T_\varphi(\MN{D})=\mathcal X_T$,
        \item there are positive semidefinite matrices $\rho_i$ with
              $\mathcal X_T=\spn\{\rho_i\}$,
        \item $T(\mathcal X_T)=\mathcal X_T$.
    \end{enumerate}
    This is \cite[Proposition~6.12 (Asymptotic image)]{Wolf2012Quantum}.
\end{theorem}

\begin{proof}\leanok
    \uses{def:peripheral_spectral_projections,
        thm:peripheral_spectral_projection_properties,
        thm:peripheral_jordan_trivial,
        cor:peripheral_projection_channel,
        cor:stationary_density_span}
    Write $V_\mu=\ker(T-\mu)$ for the eigenspace of a peripheral eigenvalue
    $\mu$, so that $\mathcal X_T=\sum_{|\mu|=1}V_\mu$ by the definition of
    $\mathcal X_T$. The image of the projection $T_\varphi$ is its set of
    fixed points, and the fixed points of $T_\varphi$ are the linear
    combinations of peripheral eigenvectors:
    \begin{align}
        T_\varphi(\MN{D})
         & =\{X\in\MN{D}\mid T_\varphi(X)=X\}
        \label{eq:channel_asymptotic_fixed_points}\\
         & =\sum_{|\mu|=1}V_\mu
        \label{eq:channel_asymptotic_eigenspace_sum}\\
         & =\mathcal X_T.
        \label{eq:channel_asymptotic_definition}
    \end{align}
    The equality (\ref{eq:channel_asymptotic_fixed_points}) holds because
    $T_\varphi$ is a projection; a matrix fixed by $T_\varphi$ lies in the
    peripheral subspace, whose Jordan blocks are trivial for a positive
    trace-preserving map (Theorem~\ref{thm:peripheral_jordan_trivial}),
    which gives (\ref{eq:channel_asymptotic_eigenspace_sum}); and
    (\ref{eq:channel_asymptotic_definition}) is the definition of
    $\mathcal X_T$. This is the first claim.

    The projection $T_\varphi$ is itself positive and trace-preserving, so its
    fixed-point space is spanned by its stationary density matrices:
    \begin{align}
        \mathcal X_T=\{X\in\MN{D}\mid T_\varphi(X)=X\}
        =\spn\{\rho\mid \rho\ \text{a stationary density matrix of }T_\varphi\}.
    \end{align}
    Every such $\rho$ is positive semidefinite, which is the second claim.

    On $V_\mu$ the map $T$ acts as multiplication by $\mu$, and $|\mu|=1$
    forces $\mu\neq0$, so multiplication by $\mu$ carries $V_\mu$ onto itself
    and $T(V_\mu)=\mu V_\mu=V_\mu$. Summing over the peripheral eigenvalues,
    \begin{align}
        T(\mathcal X_T)=T\Bigl(\sum_{|\mu|=1}V_\mu\Bigr)
        =\sum_{|\mu|=1}T(V_\mu)=\sum_{|\mu|=1}V_\mu=\mathcal X_T,
    \end{align}
    which is the third claim.
\end{proof}

\begin{theorem}[Ces\`aro fixed-point theorem]\label{thm:cesaro_fixedpoint}
    \lean{IsChannel.exists_posSemidef_fixedPoint}
    \leanok
    \uses{def:channel, def:density_matrices, def:cesaro_mean}
    Every quantum channel $E:\MN{D}\to\MN{D}$, with $D\geq1$, has a nonzero
    positive semidefinite fixed point: there exists $\rho\geq0$, $\rho\neq0$,
    with $E(\rho)=\rho$.
\end{theorem}

\begin{proof}
    \leanok
    \uses{thm:density_nonempty, lem:cesaro_mean_density,
        thm:cesaro_telescope}
    Start with any $\rho_0\in\mc{D}_D$
    (Theorem~\ref{thm:density_nonempty}).
    Lemma~\ref{lem:cesaro_mean_density} keeps every $\sigma_N$, $N\geq1$,
    in the compact set $\mc{D}_D$, so a subsequence converges to some
    $\rho\in\mc{D}_D$. The telescope identity
    (\ref{eq:channel_cesaro_telescope}) gives
    \begin{align}
        \lVert E(\sigma_N)-\sigma_N\rVert
         & =\frac{1}{N}\lVert E^N(\rho_0)-\rho_0\rVert
        \longrightarrow0,
        \notag
    \end{align}
    because the iterates $E^N(\rho_0)$ lie in the compact set
    $\mc{D}_D$ and are therefore uniformly bounded. Along the convergent
    subsequence, continuity of $E$ now gives $E(\rho)=\rho$.
    In~\cite[Theorem~6.11]{Wolf2012Quantum}, existence is proved
    via Brouwer's fixed-point theorem; we use the Ces\`aro
    argument instead.
\end{proof}
